\section{How this connects to prior work}
\label{sec:related}

\emph{Predictive runtime verification and enforcement.}  Predictive runtime
verification uses a priori knowledge of the system to reach a conclusive verdict
before the observed execution alone permits one \cite{pinisetty2017prv}, and
predictive runtime enforcement reduces to non-predictive enforcement when no
such knowledge is available \cite{pinisetty2017pre}.  For stochastic linear
systems, Yoon \emph{et al.}\ determine whether a control sequence exists over a
finite horizon that maintains a safety property with sufficiently high
probability \cite{yoon2019prm}.  RADICS combines a black-box monitor that
switches to a safe algorithm on a failing trajectory with a white-box monitor
that predicts such trajectories earlier \cite{wheatman2021radics}.  Predictive
safety filters solve an online finite-horizon problem to certify or modify a
proposed input under model uncertainty \cite{wabersich2021filter}.  I ask a
narrower comparative question: once recovery timing and a tuned current-margin
threshold are included, when does an earlier trigger change the measured
outcome?

\emph{Runtime-assurance architectures.}  Simplex uses a simple, reliable core to
bound the behavior of more complex functionality \cite{sha2001simplex}.  Its
black-box extension uses runtime checks to transfer control between advanced and
baseline controllers and includes a multi-robot case study
\cite{mehmood2022blackbox}, and SOTER composes robotics runtime-assurance modules
that switch between advanced and safe controllers \cite{desai2019soter}.
Mission-level monitors can reject commands that would make a mission infeasible
before they execute \cite{tsai2026missionrta}; compositional monitoring can
detect policy violations that emerge only across a robot swarm
\cite{kekatos2026swarmrta}; and collective run-time adaptation addresses safety
and mission completion in multi-robot systems \cite{bozhinoski2019collective}.
These works design what happens after a trigger.  I study the trigger rule
itself: a projected time to threshold versus a fixed current-margin threshold.

\emph{Adaptive margins and time-to-threshold triggers.}  Time to collision, the
ratio of range to closing speed, has long been used to time braking and to grade
traffic conflicts \cite{lee1976ttc,mahmud2017surrogate}.
\Cref{prop:reduction} is the corresponding observation for a generic monitored
margin: a threshold on projected time is a threshold on the margin scaled by the
estimated rate.  Adaptive margins can also scale a safety buffer with relative
speed, braking capability, and prediction uncertainty \cite{xu2026raam}, and
control barrier functions provide a broader framework for encoding and enforcing
safety constraints \cite{ames2019cbf}.  Because the horizon trigger studied here
is one particular rate-scaled margin, it should be compared with tuned fixed
margins and not only with an unpadded threshold.

\emph{Prognostics and recoverability.}  Prognostics estimates whether and when
faults will occur in order to inform condition-based maintenance
\cite{jardine2006prognostics}.  Runtime-assurance verification formalizes the
adequacy of decision timing and time-bounded recoverability
\cite{nigam2023recoverability}.  Recent work on adaptive governance of AI agents
predicts when a viability index will cross its boundary from a first-order fit
of its recent trend \cite{marin2026governance}.  These results motivate a
point-of-no-return analysis, but they do not by themselves establish when an
earlier trigger beats a tuned fixed margin for a robot model.
